\clearpage
\begingroup
\setlength{\parskip}{1pt}
\titlespacing*{\section}{0pt}{0.8em}{0.45em}
\titlespacing*{\subsection}{0pt}{0.45em}{0.25em}
\section*{Reading an Instruction Description}\label{page:reading-an-instruction-description}
\phantomsection
\addcontentsline{toc}{section}{Reading an Instruction Description}
\subsection*{Entry Organization}
Each instruction entry presents its normative Operation, then Assembler Syntax, Attributes, and structured status
metadata. Detailed Semantics follows when more space is needed, before the concrete encoding forms. Brief,
non-normative summaries appear only in the instruction-set summary tables. Each mnemonic begins on a new page so its
form diagrams, field explanations, and side-effect text stay together as one reference unit.

An entry's Operation and Detailed Semantics are normative and may refine only
named common rules. Determine every observable state transition or action---trap/fault,
commit, repeat, event, flag effect, or memory action---from those fields together
with named common rules, including the
\hyperref[page:instruction-execution-model]{Instruction Execution Model}.

\subsection*{Mnemonics and Suffixes}
The B, W, L, and Q suffixes select 8-, 16-, 32-, and 64-bit integer sizes; S and D select 32- and 64-bit
floating-point scalar sizes where defined.

When an instruction form encodes a size selector z, its definition maps z to the suffix set shown for that form. A
mnemonic without a size suffix has a single architecturally defined size or no data-size operand.
Conditional mnemonic forms use the condition names listed in the Flags and Condition Codes chapter.

CMPJcc, DJcc, IJcc, Jcc, MOVcc, SETcc, TESTJcc, and FMOVcc place the condition suffix in the mnemonic. A form with
only one architectural size omits the size suffix.

\subsection*{Operand and Status Notation}
Rn names a general-purpose register; SP and PC are special registers outside that encoding namespace. An
\texttt{<ea>} operand names a compact effective-address field and any appended EA payload.

FLAGS, STATUS, FFLAGS, and FSTATUS are architectural state registers. Instruction descriptions state whether a
form updates, preserves, reads, or ignores each status field.

In instruction pseudocode, Operand Read applies the addressing mode and operation size. Operand Write stores the
selected-size result; a sub-Q Rn write preserves upper bits unless the instruction defines another rule. EA Address
computes an address without reading memory. A control-target memory form reads its declared-width target, while a
control-target register or immediate form supplies the target value directly. A selector may name an Rn register, a
segment register, or an encoded immediate. FLAGS ZNCV refers to the integer Z, N, C, and V bits, while FFLAGS names
the accrued floating-point exception and status bits.

\subsection*{Encoding Forms}
In an encoding diagram, fixed 0 and 1 fields identify opcode bits, and the notes below decode lettered fields. An
<ea> field is the compact effective-address field and may append EA payload bytes in operand order.

In each instruction entry, Encoding Forms presents the concrete forms, operand and
effective-address grammar, stable form identities, and size selectors.

For an extended encoding, the four-bit L field selects a 3+L-byte instruction record. The required-byte count is the
space consumed by the opcode and selected operand payloads. Every larger encoded length through 18 bytes is
valid and the trailing bytes are uninterpreted padding. The \texttt{LEN n,} spelling records an explicit total length.
\endgroup
\clearpage
